\scanpage{23}
\mspage{5}

Ainsi, le sous-$\unsure{\mathfrak S_3|\mathcal G_3}$-ens. \inserted{$X'_3$} de $X_3$ engendré par $d(X_2)$ a une unique orbite à 1 \unsure{élémt}, savoir $\{e_3\}=\{de_2\}$, toutes les autres orbites sont de cardinal 3 \deleted{et} (avec la représentation régulière de $\unsure{\mathfrak S_3|\mathcal G_3}$), et elles sont en corr. biunivoque avec les éléments de $X_2\setminus\{e_2\}\overset{\mathrm{def}}=X_2^*$. Comme $\unsure{\mathfrak S_3|\mathcal G_3}$-ens $X'_3$ s'écrit canoniquement
\[
X'_3\simeq\{e_3\}\amalg\left(\unsure{\mathfrak S_3|\mathcal G_3}\wedge_{\{1,\sigma_1\}}X_2^*\right),
\]
où $\{1,\sigma_1\}$ opère trivialement sur $X_2^*$.

\margin{Et on a une explicitation de $p$ sur $X'_3$ via les trois premières des cinq formules $(*)$, \unsure{le reste} à satisfaire \unsure{étant} \unsure{formel} \unsure{par ces trois formules} \unsure{cas}.}

Reste à exprimer
\[
X_3^*\overset{\mathrm{def}}=X_3\setminus\underbrace{X'_3}_{\unsure{\mathfrak S_3|\mathcal G_3}\cdot d(X_2)}
\]
comme $\unsure{\mathfrak S_3|\mathcal G_3}$-ens, et l'application induite par $p$
\diagram{\begin{tikzcd}[ampersand replacement=\&]
X_3^*\arrow[r,two heads,"p^*=p\mid X_3^*"]\&X_2
\end{tikzcd}}
satisfaisant à
\[
\boxed{p^*\sigma_3=\sigma p^*}.
\]

\diagram{\begin{tikzcd}[ampersand replacement=\&]
X_3\arrow[r,two heads,"{p,p\rho}"]\&X_2\times X_2
\end{tikzcd}}
épi

\margin{$p(x,y,z)=(x,y)$
\[
p'(x,y,z)=(y,z)=p\rho(x,y,z)
\]}

pour $\unsure{\zeta|\xi},\eta\in X_2$, $\exists\alpha\in X_3$ avec
\[
p\alpha=\unsure{\zeta|\xi},\quad p\rho\alpha=\eta
\]

\margin{22}
